% \begin{abstract}
\begin{abstract}
We study the allocation of indivisible chores with subsidies under negative dichotomous valuations, where the marginal disutility of
every chore is either zero or one. This is the chore analogue of the dichotomous goods model of Barman et al.~\cite{BKNS22}, for which a subsidy of at most one unit per agent is known to suffice for envy-freeness.

We show that the same optimal guarantee holds for chores, under no structural assumption beyond binary marginals. For every instance with negative dichotomous valuations, there exists a complete allocation $\mathcal{A}$ and a subsidy vector $p\in\{0,1\}^n$ with
\(
\sum_{i\in N} p_i\le n-1
\) such that $(\mathcal{A},p)$ is envy-free. Moreover, the allocation $\mathcal{A}$ is EF1 even before subsidies are paid. The bound is
tight, and the allocation and subsidies can be computed in polynomial time in the value-oracle model.

Our algorithm starts from an envy-free partial allocation produced by the algorithm of Tao et al.~\cite{TWYZ23}, assigns the remaining chores to distinct agents in a tail strongly connected component selected at termination, and determines the subsidised agents through a backward closure in the associated equality graph.
\end{abstract}


% We study the problem of allocating indivisible \emph{chores} among agents in a
% fair manner, with subsidies. Envy-free allocations of indivisible items are not
% guaranteed to exist, but envy-freeness can be achieved by additionally providing
% some subsidy to the agents; in the chores setting th is subsidy is naturally read
% as compensation paid to the agents who absorb the workload. We study the subsidy
% required under \emph{negative dichotomous} valuations, i.e., among agents whose
% marginal disutility for any chore is either zero or one:
% $v_i(S) - v_i(S \cup \set{g}) \in \set{0,1}$ for every agent $i$, every
% $S \subseteq \items$ and every $g \in \items \setminus S$. This is the exact
% mirror of the dichotomous goods model of Barman et al.~\cite{BKNS22}, who prove
% that a per-agent subsidy of $0$ or $1$ always suffices. We prove that the
% same guarantee holds for chores.

% Existing work already gives a bounded subsidy: negative dichotomous valuations
% are doubly monotone under the two-sided normalisation, so the reduction of
% Kawase et al.~\cite{KMSTY24} yields $n-1$ per agent and $n(n-1)/2$ in total. The
% question is whether the far stronger $\set{0,1}$ guarantee of~\cite{BKNS22}
% survives the passage from goods to chores, yielding a linear-factor improvement over that baseline.
% The Halpern--Shah characterisation of envy-freeable
% allocations~\cite{HS19} transfers verbatim, since it assumes nothing about the
% sign or the structure of the valuations, and it delivers an \emph{integral}
% minimum subsidy vector here; the matching lower bound of $n-1$ transfers as
% well, so the bound in question is tight.

% We prove that bound, for every number of agents: every negative dichotomous
% instance admits an envy-free solution with a subsidy of $0$ or $1$ per agent,
% hence at most $n-1$ in total, computable in polynomial time in the value-oracle
% model. The allocation we construct is moreover EF1 before any subsidy is
% paid, so the money is what upgrades it from EF1 to exact envy-freeness. All
% of this holds under no assumption on the cost functions beyond binary marginals ---
% not additivity, not submodularity, no relation between different agents' cost
% functions, and no bound on the number of chores. The proof takes an envy-free
% \emph{partial} allocation, produced by the algorithm of Tao et
% al.~\cite{TWYZ23} and leaving at most $n-1$ chores unassigned, and completes it:
% the residual chores go to distinct agents inside the strongly connected
% component at which that algorithm halts, and the paid set is the backward
% closure of those recipients under the relation ``agent $i$ values her own bundle
% exactly as she values agent $j$'s''. Envy-freeness of the completion is verified
% directly, so the argument does not appeal to the characterisation of
% envy-freeable allocations.
% \end{abstract}